\chapter{Využití nástrojů umělé inteligence}
\label{app:ai-disclosure}

Při přípravě práce byl využit nástroj Claude (Anthropic, modely Opus~4
a~Sonnet~4) jako AI asistent v~následujících oblastech:

\begin{description}
    \item[Psaní textu práce.]
        Nástroj byl využit pro drafting a~úpravu textových částí
        (kapitoly~1--5), kontrolu konzistence argumentace a~LaTeX
        formátování. Autor formuloval záměr a~strukturu každé sekce;
        AI asistent navrhoval formulace, které autor revidoval
        a~upravoval. Výzkumný design, argumentace a~finální znění
        jsou autorovy vlastní.

    \item[Vyhledávání a~shrnutí literatury.]
        Nástroj byl využit k~identifikaci relevantních zdrojů
        a~shrnutí obsahu papers. Autor ověřoval relevanci a~správnost
        každého zdroje před zařazením do textu; citace odkazují
        výhradně na zdroje přečtené a~posouzené autorem.

    \item[Diagnostická analýza v~experimentu.]
        V~kroku Diagnóza iteračního cyklu (sekce~\ref{sec:pilotni-iterace})
        byl nástroj využit jako analytický asistent: výzkumník
        popisoval naměřené výsledky, AI asistent navrhoval možné
        příčiny selhání a~opravy instrukcí na základě literatury.
        Finální rozhodnutí o~každé změně \texttt{AGENTS.md} bylo
        vždy na autorovi.

    \item[Implementace měřicí infrastruktury.]
        Experimentální skripty (\texttt{new-run.ts},
        \texttt{analyze-run.ts}, \texttt{judge.ts}) byly
        implementovány s~asistencí AI nástroje na základě autorova
        návrhu specifikace metrik. Autor navrhl architekturu,
        definoval požadované výstupy a~provedl validaci správnosti
        měření (příloha~\ref{app:ai-disclosure}).
\end{description}

\noindent
AI nástroj nepouštěl experimenty, nevyhodnocoval výsledky
a~nerozhodoval o~směru výzkumu. Za celý obsah práce nese
plnou odpovědnost autor.
